% ===== PRÉAMBULE CANONIQUE — à coller VERBATIM en tête de CHAQUE .tex =====
\documentclass[11pt,a4paper]{article}
\usepackage[utf8]{inputenc}\usepackage[T1]{fontenc}\usepackage{lmodern}
\usepackage[french]{babel}
\usepackage{amsmath,amssymb,mathtools}\usepackage{icomma}
\usepackage{siunitx}\sisetup{locale=FR}  % \qty{}{}, \si{}, \ang{}
\usepackage{xcolor}\usepackage{colortbl}
\usepackage{tikz}\usepackage{pgfplots}\pgfplotsset{compat=1.18}
\usepackage[siunitx,europeanresistors]{circuitikz} % circuits (élec.)
\usepackage[most]{tcolorbox}\usepackage{enumitem}\usepackage{array,booktabs}
\usepackage[a4paper,margin=2.2cm]{geometry}
\usetikzlibrary{arrows.meta,calc,angles,quotes,positioning,patterns,%
  backgrounds,shapes.geometric,decorations.pathmorphing,decorations.markings,intersections}
\definecolor{primary}{RGB}{222,49,99}\definecolor{primarydark}{RGB}{150,22,55}
\definecolor{defcol}{RGB}{0,114,178}\definecolor{methcol}{RGB}{52,92,148}
\definecolor{excol}{RGB}{0,135,90}\definecolor{attncol}{RGB}{200,40,55}
\tcbset{boxbase/.style={enhanced,breakable,boxrule=0.9pt,arc=2.5pt,
  left=3mm,right=3mm,top=2.3mm,bottom=2.3mm,fonttitle=\bfseries,coltitle=white}}
% --- boîtes TUTORIEL (noms EXACTS, ne pas renommer) ---
\newtcolorbox{cle}[1][]{boxbase,colback=defcol!6,colframe=defcol,
  title={\textsc{À retenir}\ifx\relax#1\relax\relax\else\textmd{ : \itshape #1}\fi}}
\newtcolorbox{methode}[1][]{boxbase,colback=methcol!7,colframe=methcol,sharp corners,
  title={\textsc{Méthode}\ifx\relax#1\relax\relax\else\textmd{ --- \itshape #1}\fi}}
\newtcolorbox{exemplebox}[1][]{boxbase,colback=excol!5,colframe=excol,
  title={\textsc{Exemple}\ifx\relax#1\relax\relax\else\textmd{ : \itshape #1}\fi}}
\newtcolorbox{piege}[1][]{boxbase,colback=attncol!6,colframe=attncol,
  title={\textsc{Piège}\ifx\relax#1\relax\relax\else\textmd{ : \itshape #1}\fi}}
\newtcolorbox{remarque}[1][]{boxbase,colback=black!4,colframe=black!45,
  title={\textsc{Remarque}\ifx\relax#1\relax\relax\else\textmd{ : \itshape #1}\fi}}
% --- boîtes EXERCICE / CORRIGÉ (noms EXACTS, auto-comptées) ---
\newtcolorbox[auto counter]{exo}[1][]{boxbase,colback=primary!4,colframe=primary,
  title={\textsc{Exercice}~\thetcbcounter\ifx\relax#1\relax\relax\else\textmd{ --- \itshape #1}\fi}}
\newtcolorbox[auto counter]{corrige}[1][]{boxbase,colback=excol!4,colframe=excol,
  title={\textsc{Corrigé}~\thetcbcounter\ifx\relax#1\relax\relax\else\textmd{ --- \itshape #1}\fi}}
\newcommand{\vv}[1]{\vec{#1}}\newcommand{\ds}{\displaystyle}
\newcommand{\R}{\mathbb{R}}\newcommand{\N}{\mathbb{N}}\newcommand{\Z}{\mathbb{Z}}
\begin{document}
\sisetup{per-mode=symbol}

\begin{exo}[à quelle distance placer l'objet ?]
On impose la taille de l'image et on cherche la position de l'objet.
\begin{enumerate}[label=\alph*)]
  \item Une lentille \textbf{divergente} de focale $f=-\qty{20}{\cm}$ doit donner d'un objet une image
        \emph{droite, virtuelle et trois fois plus petite} ($\gamma=+\tfrac13$). À quelle distance $u$
        placer l'objet ?
  \item Une lentille \textbf{convergente} de focale $f=+\qty{10}{\cm}$ doit projeter sur un écran une
        image \emph{réelle, renversée, quatre fois plus grande} ($\gamma=-4$). Calcule $u$ puis $v$.
\end{enumerate}
\end{exo}

\begin{exo}[le problème du facteur $x$]
Un objet lumineux est placé à une distance $u=x\times f$ d'une lentille convergente de focale $f$. On
observe une image nette et réelle sur un écran situé à $v=1{,}2\times f$ de la lentille.
\begin{enumerate}[label=\alph*)]
  \item En écrivant la relation de conjugaison avec ces expressions, montre que $x$ vérifie
        $1=\dfrac{1}{x}+\dfrac{1}{1{,}2}$, et calcule $x$.
  \item L'objet est-il au-delà de $2f$ ? L'image est-elle agrandie ou réduite ? (calcule $\gamma$.)
\end{enumerate}
\end{exo}

\begin{exo}[attention aux unités]
\begin{enumerate}[label=\alph*)]
  \item Une lentille convergente a $f=\qty{0.05}{\m}$. Un objet est placé à $d=\qty{10}{\cm}$ de la
        lentille. À quelle distance $d'$ se forme l'image ? Quelle est sa nature complète ?
  \item Un objet est à l'infini devant une lentille convergente de focale $f=\qty{3}{\cm}$. Où se
        forme l'image et quelle est sa nature ?
\end{enumerate}
\end{exo}

\begin{exo}[calculer pour vérifier une affirmation]
Une lentille convergente a $f=+\qty{5}{\cm}$. Pour chaque affirmation, fais le calcul et dis si elle
est \textbf{vraie} ou \textbf{fausse}.
\begin{enumerate}[label=\alph*)]
  \item Objet à $u=+\qty{15}{\cm}$ : « l'image est réelle, renversée, à \qty{7.5}{\cm} de la lentille ».
  \item Objet à $u=+\qty{3}{\cm}$ : « l'image est réelle et agrandie ».
  \item Objet à $u=+\qty{10}{\cm}$ : « l'image est de même taille, renversée, à \qty{10}{\cm} ».
\end{enumerate}
\end{exo}

\end{document}
